% PAPER FOR ICAPS 2008
% Considers new neighbourhood structure for fuzzy job shop (extension of Van Laarhoven, based on new def of critical paths for fuzzy disjunctive graph).
% Contribution: definition of critical paths and consequent neighbourhood, theoretical results: improvement w.r.t. previous neighbourhood, feasibility and connectivity.
% Experimental results using MA from Fuzz-IEEE 2007: comparison with previous approaches, comparison H vs. H' from Fuzz-IEEE2007, synergy

\documentclass[letterpaper]{article}

% packages
\usepackage{aaai}% AAAI formatting instructions
\usepackage{times}
\usepackage{helvet}
\usepackage{courier}

%\usepackage[dvips]{graphicx}
\usepackage[dvips]{graphicx}
%\usepackage[hang]{subfigure}% to include subfigures in a figure environment
\usepackage{amsmath, amssymb, amsthm}% to use amsmath typesetting options


% theorem environments
\newtheorem{definition}{Definition}
\newtheorem{theorem}{Theorem}
\newtheorem{proposition}{Proposition}
\newtheorem{corollary}{Corollary}

\DeclareMathOperator{\N}{N}
\DeclareMathOperator{\Cr}{Cr}


\begin{document}

\title{A New Local Search for the Job Shop Problem with Uncertain Durations}

\author{In\'es Gonz\'alez-Rodr\'{\i}guez\\
    Dept. of Mathematics, Statistics and Computing, University of Cantabria (Spain)\\ \texttt{ines.gonzalez@unican.es}
\AND Camino R. Vela \and Jorge Puente \and Ramiro Varela \\
    Dept. of Computer Science and A.I. Centre, University of Oviedo (Spain)\\ \texttt{\{crvela,puente,ramiro\}@uniovi.es}
}

\maketitle

\begin{abstract}
In the sequel we consider the job shop scheduling problem with uncertain durations represented as triangular fuzzy numbers. We propose a new neighbourhood structure for local search, based on a new definition of critical path for fuzzy durations. A theoretical analysis of the proposed structure shows that it improves a previous one from the literature. It also shows that feasibility and connectivity hold, these being two highly desirable properties. Experimental results are reported which further illustrate the potential of the proposal.%82 words
\end{abstract}

%------------------------------------------------------
% INTRODUCTION
%------------------------------------------------------
%
\section{Introduction}\label{secIntro}
%Introduction: motivation, previous work
Scheduling problems form an important body of research since the late fifties, with multiple applications in industry, finance and science~\cite{Brucker2006}. Part of this research is devoted to fuzzy scheduling, in an attempt to model the uncertainty and vagueness pervading real-world situations. The approaches are diverse, ranging from representing incomplete or vague states of information to using fuzzy priority rules with linguistic qualifiers or preference modelling~\cite{Dubois2003},\cite{Slowinski2000}.

Incorporating uncertainty to scheduling usually requires a significant reformulation of the problem and solving methods, in order that the problem can be precisely stated and solved efficiently and effectively. For instance, extending critical path analysis to ill-known processing times is far from being trivial~\cite{Dubois2003}. Furthermore, in classical scheduling the complexity of problems such as shop problems means that practical approaches to solving them usually involve heuristic strategies: simulated annealing, genetic algorithms, local search, etc~\cite{Brucker2006}. Some attempts have been made to extend these heuristic methods to fuzzy scheduling problems where uncertain durations are modelled via fuzzy intervals, most commonly and successfully for the flow shop problem: among others, a genetic algorithm is used in \cite{Celano2003} and a genetic algorithm is hybridised with a local search procedure in \cite{Ishibuchi1998}. For the job shop, single objective problems are tackled using simulated annealing in \cite{Fortemps1997} and a genetic algorithm hybridised with local search in \cite{GonzalezRodriguez2007b}, whilst multiobjective problems are solved using genetic algorithms in \cite{Sakawa2000} or \cite{GonzalezRodriguez2008} and using a neural approach in \cite{Tavakkoli-Moghaddam2008}.

In the following, we consider a job shop problem with task durations modelled as triangular fuzzy numbers. We propose new definitions of critical path and neighbourhood structure for local search. We give theoretical results which demonstrate the potential of the new proposal, further illustrated by the experimental results.

%------------------------------------------------------
% UNCERTAINTY MODELLING
%------------------------------------------------------
%
\section{Uncertain Durations}\label{secUncertainTimes}
% - Motivation
In real-life applications, it is often the case that the exact duration of a task, i.e. the time it takes to be processed, is not known in advance. However, based on previous experience, an expert may have some knowledge (albeit uncertain) about the duration. In the literature, it is common to use fuzzy intervals to represent such uncertainty, as an alternative to probability distributions, which require a deeper knowledge of the problem and usually yield a complex calculus.

% - TFN
When there is little knowledge available, the crudest representation for uncertain processing times would be a human-originated confidence interval. If some values appear to be more plausible than others, a natural extension is a fuzzy interval or fuzzy number. The simplest model is a \emph{triangular fuzzy number} or \emph{TFN}, using only an interval $[a^1,a^3]$ of possible values and a modal value $a^2$ in it. For a TFN $A$, denoted $A=(a^1, a^2, a^3)$, the membership function takes the following triangular shape:
\begin{equation}\label{eqTFNmembership}
\mu_A(x)=\begin{cases}

\frac{x-a^1}{a^2-a^1}&: a^1 \leq x \leq a^2\\
\frac{x-a^3}{a^2-a^3}&: a^2 < x \leq a^3\\
0 &: x<a^1 \text{ or } a^3 < x\\
\end{cases}
\end{equation}

% - FUZZY INTERVAL ARITHMETIC
% Intro, preliminary concepts
Triangular fuzzy numbers and more generally fuzzy intervals have been extensively studied in the literature (cf.~\cite{Dubois1988}). A \emph{fuzzy interval} $Q$ is a fuzzy quantity (a fuzzy set on the reals) whose $\alpha$-cuts $Q_\alpha=\{r \in \mathbb{R}: \mu_Q(r) \geq \alpha\}$, $\alpha \in (0,1]$, are intervals (bounded or not). The \emph{support} of $Q$ is $Q_0=\{r \in \mathbb{R}: \mu_Q(r)>0\}$. A \emph{fuzzy number} is a fuzzy quantity whose $\alpha$-cuts are closed intervals, with compact support and unique modal value.

% arithmetic operations with fuzzy intervals: approximation
In the job shop, we essentially need two operations on fuzzy quantities, the sum and the maximum. These are obtained by extending the corresponding operations on real numbers using the \emph{Extension Principle}. However, computing the resulting expression is cumbersome, if not intractable. For the sake of simplicity and tractability of numerical calculations, we follow~\cite{Fortemps1997} and approximate the results of these operations by a TFN, evaluating only the operation on the three defining points of each TFN. The approximated sum coincides with the sum of TFNs as defined by the Extension Principle, so for any pair of TFNs $M$ and $N$:
\begin{equation}
M+N=(m^1+n^1, m^2+n^2, m^3+n^3)
\end{equation}
The same is not always true for the maximum $\vee$, although it is possible to prove that for any two TFNs $M,N$, if $F=N \vee M$ denotes their maximum and $G=(m^1\vee n^1, m^2 \vee n^2, m^3 \vee n^3)$ its approximated value, it holds that:
\begin{equation}
\forall \alpha \in [0,1], \quad \underline{f}_\alpha \leq \underline{g}_\alpha, \overline{f}_\alpha \leq \overline{g}_\alpha.
\end{equation}
where $[\underline{f}_\alpha, \overline{f}_\alpha]$ is the $\alpha$-cut of $F$. In particular, $F$ and $G$ have identical support and modal value, that is, $F_0=G_0$ and $F_1=G_1$. This approximation can be trivially extended to the case of more than two TFNs.


% - EXPECTED VALUE
The membership function $\mu_Q$ of a fuzzy quantity $Q$ can be interpreted as a possibility distribution on the real numbers, so we naturally obtain the \emph{possibility} $\Pi$ and \emph{necessity measure} $\N$ that $Q \leq r$, where $r$ is a real number~\cite{Dubois1988}. These measures are used to define the \emph{expected value} of a fuzzy quantity~\cite{Liu2002}, given for a TFN $A$ by
\begin{equation}
E[A]=\frac{1}{4}(a^1+2a^2+a^3).
\end{equation}
The expected value coincides with the the \emph{neutral scalar substitute} of a fuzzy interval and can also be obtained as the centre of gravity of its mean value or using the area compensation method~\cite{Dubois2003}. It induces a total ordering $\leq_E$ in the set of fuzzy intervals~\cite{Fortemps1997}, where for any two fuzzy intervals $M, N$ $M \leq_E N$ if and only if $E[M] \leq E[N]$. Clearly, for any two TFNs $A$ and $B$, if $\forall i, a^i \leq b^i$, then $A \leq_E B$.

%------------------------------------------------------
% JOB SHOP
%------------------------------------------------------
%
\section{The Job Shop Scheduling Problem}\label{secJobShop}
% - JOB SHOP DESCRIPTION
The classical \emph{job shop scheduling problem}, also denoted \emph{JSP}, consists in scheduling a set of jobs $\{J_1,\dots,J_n\}$ on a set $\{M_1, \dots, M_m\}$ of physical resources or machines, subject to a set of constraints. There are \emph{precedence constraints}, so each job $J_i$, $i=1,\dots,n$, consists of $m$ tasks $\{\theta_{i1},\dots,\theta_{im}\}$ to be sequentially scheduled. Also, there are \emph{capacity constraints}, whereby each task $\theta_{ij}$ requires the uninterrupted and exclusive use of one of the machines for its whole processing time. A solution to this problem is a schedule, i.e. an allocation of starting times for each task, which, besides being \emph{feasible} (i.e. all constraints hold), is \emph{optimal} according to some criterion, most commonly that the \emph{makespan} is minimal.

% - EXPECTED MAKESPAN MODEL FOR FUZZY JOB SHOP: SCHEDULE BUILDER AND PROBLEM FORMULATION
\subsection{Expected Makespan Model}\label{subsecExpectedMakespanProblem}
% SCHEDULE BUILDER
For a job shop problem instance of size $n \times m$ ($n$ jobs and $m$ machines), let $\boldsymbol{p}$ be a duration matrix and let $\boldsymbol{\nu}$ be a machine matrix such that $p_{ij}$ is the processing time of task $\theta_{ij}$ and $\nu_{ij}$ is the machine required by $\theta_{ij}$, $i=1,\dots,n$, $j=1,\dots,m$. Let $\sigma$ be a feasible task processing order, i.e., a lineal ordering of tasks which is compatible with a processing order of tasks that may be carried out so that all constraints hold. A feasible schedule may be derived from $\sigma$ using a \emph{semi-active schedule builder}. Let $C_i(\sigma,\boldsymbol{p},\boldsymbol{\nu})$ denote the completion time of job $J_i$ and let $S_{ij}(\sigma,\boldsymbol{p},\boldsymbol{\nu})$ and $C_{ij}(\sigma,\boldsymbol{p},\boldsymbol{\nu})$ denote the starting and completion times of task $\theta_{ij}$, $i=1,\dots,n$ $j=1,\dots,m$. These times are obtained as follows:
\begin{gather}
C_i(\sigma,\boldsymbol{p},\boldsymbol{\nu}) = C_{im}(\sigma,\boldsymbol{p},\boldsymbol{\nu})\\
S_{ij}(\sigma,\boldsymbol{p},\boldsymbol{\nu}) = C_{i(j-1)}(\sigma,\boldsymbol{p},\boldsymbol{\nu}) \vee
    C_{rs}(\sigma,\boldsymbol{p},\boldsymbol{\nu})\\
C_{ij}(\sigma,\boldsymbol{p},\boldsymbol{\nu}) = S_{ij}(\sigma,\boldsymbol{p},\boldsymbol{\nu}) + p_{ij}
\end{gather}
where $\theta_{rs}$ is the task preceding $\theta_{ij}$ in the machine according to the processing order $\sigma$, $C_{i0}(\sigma,\boldsymbol{p},\boldsymbol{\nu})$ is assumed to be zero and, analogously, $C_{rs}(\sigma,\boldsymbol{p},\boldsymbol{\nu})$ is taken to be zero if $\theta_{ij}$ is the first task to be processed in the corresponding machine. The \emph{makespan} $C_{max}(\sigma, \boldsymbol{p}, \boldsymbol{\nu})$ is the maximum completion time of jobs $J_1,\dots,J_n$:
\begin{equation}
C_{max}(\sigma, \boldsymbol{p}, \boldsymbol{\nu})= \vee_{1 \leq i \leq n}\left( C_i(\sigma,\boldsymbol{p},\boldsymbol{\nu}) \right)
\end{equation}
% Notation
For the sake of a simpler notation, we may write, for instance, $C_{max}(\sigma)$ when the problem (hence $\boldsymbol{p}$ and $\boldsymbol{\nu}$) is fixed or even $C_{max}$ when no confusion is possible regarding the task processing order.

% Remark for fuzzy case
If task processing times are fuzzy intervals, then the addition and maximum in the equations above are taken to be the corresponding operations on fuzzy intervals, approximated for the particular case of TFNs as proposed above. The obtained schedule will be a fuzzy schedule, where the starting and completion times of all tasks and the makespan are fuzzy intervals. Such intervals may be seen as possibility distributions on the values that the times may take. However, the task processing ordering $\sigma$ that determines the schedule is crisp; there is no uncertainty regarding the order in which tasks are to be processed. In other words, we obtain a fuzzy schedule from a crisp task ordering.

% Example
To illustrate the schedule builder, consider a problem of 3 jobs and 2 machines with the following matrices for fuzzy processing times and machine allocation:
\begin{equation*}
\boldsymbol{p}=\begin{pmatrix}
    ( 3, 4, 7) & ( 1, 2, 3)\\
    ( 4, 5, 6) & ( 2, 3, 4)\\
    ( 1, 2, 6) & ( 1, 2, 4)
  \end{pmatrix}
  \boldsymbol{\nu}=\begin{pmatrix}
    1 & 2\\
    2 & 1\\
    2 & 1
  \end{pmatrix}
\end{equation*}
Figure~\ref{figSchedule} shows the Gantt chart (adapted to TFNs following \cite{Fortemps1997}) of the schedule given by the task order $\theta_{11}, \theta_{21}, \theta_{31},\theta_{22}, \theta_{12},\theta_{32}$. For instance the starting and completion times for task $\theta_{22}$ will be $S_{22}=C_{21} \vee C_{11}=(4,5,7)$ and $C_{22}=S_{22}+ p_{22}=(6,8,11)$.
\begin{figure}[htb]
\centerline{\includegraphics[width=0.90\columnwidth]{schedule3.eps}}
\caption{\label{figSchedule} Gantt chart of the schedule represented by $\theta_{11}, \theta_{21}, \theta_{31},\theta_{22}, \theta_{12},\theta_{32}$}%the decision variable (1 2 3 2 3 1)}
\end{figure}

% EXPECTED MAKESPAN MODEL
% Definition of "optimal" schedule
Since we may build a feasible schedule from a feasible task processing order, we restate the goal of the job shop problem as that of finding an optimal task processing order, in the sense that makespan for the derived schedule is optimal. However, it is not trivial to optimise a fuzzy makespan, since neither the maximum nor its approximation define a total ordering in the set of TFNs. In the literature, this problem is tackled using ranking methods, lexicographical orderings, comparisons based on $\alpha$-cuts or defuzzification methods. Here we use the concept of expected value for a fuzzy quantity and the total ordering it provides, and consider that the objective is to minimise the expected makespan $E[C_{max}(\sigma, \boldsymbol{p}, \boldsymbol{\nu})]$, a crisp objective function. An optimal task processing order is then defined as follows:
\begin{definition}
Let $\Sigma$ be the set of feasible task processing orders. A task processing order $\sigma_0 \in \Sigma$ is said to be \emph{optimal} if and only if $E[C_{max}(\sigma_0)] = \min{ \{E[C_{max}(\sigma)]: \sigma \in \Sigma \}}$
\end{definition}

% Problem formulation (expression) using decision variable
A feasible task processing order may be represented by a decision variable $\boldsymbol{z}=(z_1,\dots,z_{nm})$, where $1 \leq z_l \leq n$ for $l=1,\dots,nm$ and $|\{z_l: z_l=i\}|=m$ for $i=1,\dots,n$. This is a permutation with repetition~\cite{Bierwirth1995}, a permutation of the set of tasks, where each task is represented by the number of its job, so the order of precedence among tasks requiring the same machine is given by the order in which they appear in the decision variable $\boldsymbol{z}$. For instance, the processing order from the example in Figure~\ref{figSchedule} is represented by the decision variable (1 2 3 2 1 3). Using this representation, the \emph{expected makespan model} for the job shop~\cite{GonzalezRodriguez2007b} is expressed as follows:
\begin{equation}\label{eqExpMakespanModel}
\begin{cases}
\min & E[C_{max}(\boldsymbol{z}, \boldsymbol{p},\boldsymbol{\nu})]\\
\text{subject to:} & \\
    & 1 \leq z_l \leq n, \quad l=1,\dots,nm,\\
    & |\{z_l: z_l=i\}|=m, \quad i=1,\dots,n,\\
    & z_l \in \mathbb{Z}^+, \quad l=1,\dots,nm.
\end{cases}
\end{equation}

% - DISJUCTIVE GRAPH MODEL REPRESENTATION
\subsection{The Disjunctive Graph Model Representation}\label{subsecGraphModel}

% Disjunctive graph
A job shop problem instance may be represented by a directed graph $G = (V, A \cup D)$. Each node in the set $V$ represents a task of the problem, with the exception of the dummy nodes \emph{start} or 0 and \emph{end} or $nm+1$, representing tasks with null processing times. Task $\theta_{ij}$, $1\leq i \leq n$, $1\leq j \leq m$, is represented by node $x=m(i-1)+j$. Arcs in $A$ are called \emph{conjunctive arcs} and represent precedence constraints (including arcs from node \emph{start} to the first task of each job and arcs form the last task of each job to the node \emph{end}). Arcs in $D$ are called \emph{disjunctive arcs} and represent capacity constraints. Set $D$ is partitioned into subsets $D_i$, $D=\cup_{i=1,\dots,m} D_i$, where $D_i$ corresponds to machine $M_i$ and includes an arc for each pair of tasks requiring that machine. Each arc is weighted with the processing time of the task at the source node (a TFN in our case). A feasible schedule task processing order $\sigma$ is represented by an acyclic subgraph $G(\sigma)$ of $G$, $G(\sigma)=(V,A \cup R(\sigma))$, where $R(\sigma)=\cup_{i=1 \dots m}R_i(\sigma)$, $R_i(\sigma)$ being a hamiltonian selection of $D_i$.
% figure, example
Figure~\ref{figFuzzyDisyuntiveGraph} shows the solution graph for the example in Figure~\ref{figSchedule}.
\begin{figure}[htb]
\centerline{\includegraphics[width=0.75\columnwidth]{grafo-disyuntivo-TFN.eps}}
\caption{\label{figFuzzyDisyuntiveGraph} Solution graph $G(\sigma)$ for the processing order $\sigma=$(1 3 5 4 2 6) corresponding to the Gantt chart of Figure~\ref{figSchedule}. $C_{max}(\sigma) = (7, 10, 16)$}
\end{figure}

% Critical paths in the solution graph
In the crisp case, a \emph{critical path} is defined as the longest path from node \emph{start} to node \emph{end} and a \emph{critical arc} or \emph{critical activity} is an arc or activity in a critical path. It is not trivial to extend these concepts and related algorithms to the problem with uncertain durations (cf.~\cite{Dubois2003}). For the fuzzy job shop considered herein it may even be the case that the makespan (a TFN) does not coincide with the completion time of one job (unlike the crisp case).

% - LOCAL SEARCH
\subsection{Local Search}\label{subsecLocalSearch}
% Local search schema based on critical arcs
Part of the interest of critical paths stems from the fact that they may be used to define neighbourhood structures for local search. Roughly speaking, a typical local search schema starts from a given processing order, calculates its neighbourhood and then neighbours are evaluated (for instance, using the aforementioned semiactive schedule builder) in the search of an improving solution. In \emph{steepest descent hill-climbing}, all neighbours are evaluated, the best one is selected and it replaces the original solution if it is an improving neighbour (in our case, with a smaller $E[C_{max}]$). In \emph{simple hill-climbing}, evaluation stops as soon as the first improving neighbour is found. Local search starts again from that improving neighbour, so the procedure finishes when no neighbour satisfies the acceptation criterion.

% neighbourhood of VanLaarhoven1992
Clearly, a central element in any local search procedure is the definition of neighbourhood. For the crisp job shop, a well-known neighbourhood, which relies on the concepts of critical path and critical block, is that proposed in~\cite{VanLaarhoven1992}.
\begin{definition}
A \emph{move} for a feasible task processing order $\sigma \in \Sigma$ is defined as the change in the order of a pair of consecutive tasks $(x,y)$ in a critical block. The \emph{neighbourhood} of $\sigma$ is defined as the set of processing orders obtained from $\sigma$ after applying all possible moves.
\end{definition}
This neighbourhood presents certain interesting characteristics in the crisp case. First, all neighbours considered represent feasible solutions and the reversal of a non-critical arc in the graph can never reduce the makespan. Therefore, it reduces the search space without any loss in the potential quality of solutions. Also, there exists a finite sequence of transitions leading from any given element to some globally optimal element. This is usually referred to as connectivity property in the literature and, as we shall argue in the sequel, is a desirable (but not general) property of neighbourhood structures.

%------------------------------------------------------
% NEW NEIGHBOURHOOD FOR LOCAL SEARCH
%------------------------------------------------------
%
%\section{The Parallel Disjunctive Graph Model Representation}\label{secParallelGraphModel}
\section{New Criticality Model and Neighbourhood Structure}
We now propose a new neighbourhood structure for the job shop with uncertain durations represented as TFNs. To do so, we give a new definition of criticality, based on the fact that all arithmetic operations used in the scheduling process are performed on the three defining points or components of the TFNs. This allows to decompose the solution graph with fuzzy durations in three parallel graphs with crisp durations.
% - PARALLEL DISJUNCTIVE GRAPH
\begin{definition}
Let $\sigma \in \Sigma$ be a feasible task processing order and let $G(\sigma)=(V,A \cup R(\sigma))$ be the corresponding solution graph, where the cost of any arc $(x,y) \in A \cup  R(\sigma)$ is a TFN representing the processing time $p_x$ of task $x$. From $G(\sigma)$, we define the \emph{parallel solution graphs} $G^i(\sigma)$, $i=1,2,3$, which are identical to $G(\sigma)$ except for the cost of arc $(x, y)\in A \cup  R(\sigma)$, which for graph $G^i(\sigma)$ will be the $i$-th defining point of $p_x$, that is, $p_x^i$.
\end{definition}
%% FIGURE WITH 3 PARALLEL SOLUTION GRAPHS
\begin{figure}[tb]
\begin{minipage}{0.99\columnwidth}
\centering \includegraphics[width=0.75\columnwidth]{grafo-disyuntivo-1.eps}\\
\centerline{$G^1$. Critical path = (3 4 6). $C_{max}^1 = 7$}
\end{minipage}\\
\begin{minipage}{0.99\columnwidth}
\centering \includegraphics[width=0.75\columnwidth]{grafo-disyuntivo-2.eps}\\
\centerline{$G^2$. Critical path = (3 4 6). $C_{max}^2 = 10$}
\end{minipage}\\
\begin{minipage}{0.99\columnwidth}
\centering \includegraphics[width=0.75\columnwidth]{grafo-disyuntivo-3.eps}\\
\centerline{$G^3$. Critical path = (3 5 6). $C_{max}^3 = 16$}
\end{minipage}
\caption{\label{figParallelDisjGraph} Parallel graphs corresponding to the graph in Figure~\ref{figFuzzyDisyuntiveGraph}, with critical paths in blue}
\end{figure}
Each of the parallel graphs $G^i(\sigma)$ is a solution graph identical to those for crisp JSP. Therefore, in each of them a critical path is the longest path from node \emph{start} to node \emph{end}. Notice that it is not necessarily unique; for instance, Figure~\ref{figParallelDisjGraph}, shows the three parallel graphs generated from the graph in Figure~\ref{figFuzzyDisyuntiveGraph}; we see that if $p_6^3 = 3$ there would be two critical paths in $G^3(\sigma)$: (3 4 6) and (3 5 6).

% - CRITICAL PATH
% Definition
Using the parallel graph representation, we may define criticality for the fuzzy job shop as follows:
\begin{definition}\label{defCritical2}
A path $P$ in $G(\sigma)$ is a \emph{critical path} if and only if $P$ is critical in some $G^i(\sigma)$. Nodes and arcs in a critical path are termed \emph{critical}. A critical path is naturally decomposed into critical blocks $B_1,\dots,B_r$, where a \emph{critical block} is a maximal subsequence of tasks of a critical path requiring the same machine. Hence, the sets of critical paths, arcs, tasks and blocks %for the FJSP
are respectively the union of critical paths, arcs, tasks and blocks in the parallel solution graphs.
\end{definition}
Unlike for the crisp JSP, we may not state that the makespan of the schedule is the cost of a critical path. However, it holds that each component of the makespan $C_{max}^i(\sigma)$ is the cost of a critical path in the corresponding solution parallel graph $G^i(\sigma)$.

% FIRST SET OF PROPERTIES (FOLLOW DIRECTLY FROM DEFINITION)
% notation: head and tails
For a solution graph $G(\sigma)$ and a task $x$, let $S_x$ and $C_x$ denote respectively the starting and completion times of $x$, let $P\nu_x$ and $S\nu_x$ denote the predecessor and successor nodes of $x$ on the machine sequence (predecessor and successor in $R(\sigma)$), and let $PJ_x$ and $SJ_x$ denote respectively the predecessor and successor nodes of $x$ on the job sequence (predecessor and successor in $A$). We now extend some well-known notions for the crisp job shop and define the \emph{head} $r_x$ of task $x$ as the starting time of $x$ i.e. $S_x$, in our context, a TFN $r_x=(r_x^1, r_x^2, r_x^3)$ given by:
\begin{equation}\label{eqHeadrx}
r_x = (r_{PJ_x} + p_{PJ_x}) \vee (r_{P\nu_x} + p_{P\nu_x}),
\end{equation}
and the \emph{tail} $q_x$ of task $x$ as:
\begin{equation}\label{eqTailqx}
q_x = (q_{SJ_x} + p_{SJ_x}) \vee (q_{S\nu_x} + p_{S\nu_x}),
\end{equation}
% elementary properties of head and tails
Clearly, $C_x = r_x + p_x$. For each parallel graph $G^i(\sigma)$, $r_x^i$ is the length of the longest path from node \emph{start} to node $x$ and $q_x^i+p_x^i$ is the length of the longest path from node $x$ to node \emph{end}. Hence, $r_x^i+p_x^i+q_x^i$ is the length of the longest path from node \emph{start} to node \emph{end} through node $x$ in $G^i(\sigma)$; it is a lower bound of the $i$-th component of the makespan, being equal to $C_{max}^i(\sigma)$ if node $x$ belongs to a critical path in $G^i(\sigma)$. For two TFNs $A$ and $B$ let $A \simeq  B$ denote that $\exists i, a^i = b^i$. The next proposition states some properties of critical arcs and tasks which follow trivially from the above definition:
% Elementary properties of critical paths (re. heads and tails)
\begin{proposition}\label{propPropertiesCriticalArcs}
If an arc $(x,y)$ is critical, then $r_x + p_x \simeq r_y$. A task $x$ is critical if and only if $r_x + p_x + q_x \simeq C_{max}$.
\end{proposition}

% - NEIGHBOURHOOD
% Definition
For the proposed definition of criticality, we may extend the neighbourhood structure defined in \cite{VanLaarhoven1992} to the fuzzy case as follows:
\begin{definition}\label{defNeighbourhoodHfuzzy}
Given an arc $v= (x,y) \in R(\sigma)$, let $\sigma_{(v)}$ denote the processing order obtained from $\sigma$ after an exchange in the processing order of tasks in arc $v$. Then, the \emph{neighbourhood structure} obtained from $\sigma$ is given by $H(\sigma) = \{ \sigma_{(v)}: v \in R(\sigma) \text{ is critical} \}$.
\end{definition}

% -- COMPARISON TO PREVIOUS APPROACHES FROM THE LITERATURE
\subsection{Comparison to Previous Approaches}\label{subsecComparison2Previous}
% Fortemps
In \cite{Fortemps1997}, we find a first proposal to extend the neighbourhood structure in \cite{VanLaarhoven1992} to the fuzzy case. An arc $(x,y)$ in the solution graph for the job shop problem with fuzzy durations is taken to be a \emph{critical arc} if and only if $C_x \simeq S_y$, that is, if it is critical in any of its components. This definition, being quite natural, yields some counterintuitive results. For instance, in the graph from %Figure~\ref{figFuzzyDisyuntiveGraph}, arc $(5,2)$ would be considered to be critical.

% GonzalezRodriguez
In ~\cite{GonzalezRodriguez2007b}, this definition was refined by incorporating backpropagation: a \emph{critical path} is taken to be a path in the solution graph from node $0$ to node $z$ where $C_z \simeq  C_{max}$ and such that for all arcs $(x,y)$ in the path it holds $C_x \simeq  S_y$. A \emph{critical arc} (respectively a \emph{critical task}) is an arc (respectively task) belonging to a critical path.

Let us denote as critical' an arc which is critical according to this last definition. All arcs considered to be critical according to \cite{Fortemps1997} will still be critical', but the reverse does not hold, avoiding some of the counterintuitive examples. For instance, in %Figure~\ref{figFuzzyDisyuntiveGraph} the critical' arcs are (1,4), (3,4), (3,5), (4,6) y (5,6) and the critical' tasks are 1, 3, 4, 5 y 6. Clearly, it is possible to use the criticality definition from \cite{GonzalezRodriguez2007b} to extend the neighbourhood structure proposed in~\cite{VanLaarhoven1992} to the fuzzy case; for a given feasible task processing order $\sigma$, let $H'(\sigma) = \{ \sigma_{(v)}: v \in R(\sigma) \text{ is critical'} \}$ denote the resulting neighbourhood.
% Proposition: comparison criticality and (consequence) H vs H'
The next proposition follows trivially from the above definitions:
\begin{proposition}\label{propComparisonCorH2InH1}
The set of critical paths is a subset of the critical' ones. Consequently, neighbourhood $H$ is contained in $H'$, that is
\begin{equation}
\forall \sigma \in \Sigma, H(\sigma) \subseteq H'(\sigma)
\end{equation}
\end{proposition}
% remarks on prop
Notice that the set of critical paths may be a proper subset of the set of critical' paths. Moreover, the second condition in Proposition~\ref{propPropertiesCriticalArcs} may not hold for a critical' task. For instance, in the example from %Figure~\ref{figFuzzyDisyuntiveGraph}, arc $(1,4)$ is critical' but is not critical and task 1 is critical' but $\forall i, r_1^i + p_1^i + q_1^i \neq C_{max}^i$.

% Proposition: neighbours outside H are non-improving neighbours
Proposition~\ref{propComparisonCorH2InH1} means a reduction in the number of neighbours when $H$ is considered instead of $H'$. We now show that this reduction does not affect the potential quality of solutions found by the local search, as the reversal of a non critical arc can never reduce the makespan:
%% PROPOSITION: REVERSAL OF A NON-CRITICAL ARC CAN NEVER REDUCE MAKESPAN
\begin{proposition}\label{propOutsideH2WorseCmax}
Let $\sigma \in \Sigma$ be a feasible processing order and let $\tau \in \Sigma$ be a feasible processing order obtained by the reversal an arc which is not critical in $G(\sigma)$. Then%The makespan for this processing order can never improve the makespan for $\sigma$:
\begin{equation}
\forall i, \quad C_{max}^i(\sigma) \leq C_{max}^i(\tau)
\end{equation}
and therefore
\begin{equation}
C_{max}(\sigma) \leq_E C_{max}(\tau)
\end{equation}
\end{proposition}
%% Proof
%\begin{proof}
\emph{Proof}Let $\tau$ be obtained from $\sigma$ by the reversal of an arc $v=(x,y)$ where $v$ is not a critical arc, that is, it is not a critical arc in any of the parallel graphs $G^i(\sigma)$. Clearly, for all $i=1,2,3$ the arcs inside critical paths of $G^i(\sigma)$ remain unchanged in $G^i(\tau)$ after the move, and therefore $C_{max}^i(\sigma) \leq C_{max}^i(\tau)$ for all $i$. \qed
%\end{proof}
%% END PROPOSITION

%% COROLLARY: H'-H DOES NOT CONTAIN ANY IMPROVING NEIGHBOURS
In particular, neighbours in $H'$ which do not belong to $H$ can never improve the makespan:
\begin{corollary}\label{propH2-H1WorseCmax}
$\forall \sigma \in \Sigma, \forall \tau \in \left( H'(\sigma) - H(\sigma) \right) \bigcap \Sigma$, $C_{max}(\sigma) \leq_E C_{max}(\tau)$
\end{corollary}

%% CONCLUSION: New neighbourhood H is ``better'' than previous neighbourhoods.
We may conclude that using $H$ instead of $H'$ in the local search procedure leads to a reduction in the number of evaluated neighbours without missing any improving neighbours, i.e. with no loss in the quality of the solution obtained by the local search. Notice that if $H'$ were to denote the neighbourhood proposed in \cite{Fortemps1997}, all the above results would still hold.

% -- MORE PROPERTIES OF NEW NEIGHBOURHOOD: THEORETICAL ANALYSIS OF FEASIBILITY AND CONNECTIVITY
\subsection{Some Desirable Properties}
Now we further study the new neighbourhood $H$, in order to prove that it has two highly desirable properties: feasibility and connectivity.
%%
%% 1) FEASIBILITY
%%
%% THEOREM: FEASIBILITY
\begin{theorem}\label{thFeasibilityH2}
Let $\sigma \in \Sigma$ be a feasible task processing order; the reversal of a critical arc $v=(x,y) \in R(\sigma)$ produces a feasible processing order, i.e., $\sigma_{(v)} \in \Sigma$. In consequence, $H(\sigma) \subset \Sigma$.
\end{theorem}
%% Proof:
%\begin{proof}
%If $v=(x,y)$ is critical, by definition,
%\begin{equation*}
%\exists i, r_y^i = r_x^i + p_x^i.
%\end{equation*}
%Suppose by contradiction that $G(\sigma_{(v)})$ has a cycle. Since $G(\sigma)$ has no cycles and the only change in $\sigma_{(v)}$ w.r.t. $\sigma$ has been the processing order between $x$ and $y$, having a cycle in $G(\sigma_{(v)})$ means that there must exist an alternative path from $x$ to $y$ in %$G(\sigma_{(v)})$, as shown in Figure~\ref{figFeasibility}.
%$G(\sigma)$, as shown in Figure~\ref{figFeasibility}.
%%Now, remember that for any task $x$, a direct consequence of the definition of the head $r_x$ in \eqref{eqHeadrxComponents} is that:
%%\begin{align}
%%\forall x, \forall i, & r_x^i \leq r_{PJ_x}^i + p_{PJ_x}^i\\
%%\forall x, \forall i, & r_x^i \leq r_{P\nu_x}^i + p_{P\nu_x}^i
%%\end{align}i
%From the definition of the head $r_y$ in \eqref{eqHeadrxComponents}, $\forall i$ we have that
%\begin{equation*}
%r_y^i \geq r_{PJ_y}^i + p_{PJ_y}^i
%\end{equation*}
%and together with the existence of an alternative path from $x$ to $y$ this means that:
%\begin{equation*}
%r_y^i \geq r_{SJ_x}^i + p_{SJ_x}^i + p_{PJ_y}^i \geq r_x^i + p_x^i + p_{SJ_x}^i + p_{PJ_y}^i
%\end{equation*}
%But being critical, there is at least one component $k$ where $r_y^k = r_x^k + p_x^k$, and for this $k$ we have:
%\begin{equation*}
%r_x^k + p_x^k \geq r_x^k + p_x^k + p_{SJ_x}^k + p_{PJ_y}^k
%\end{equation*}
%which is absurd, because all task durations are strictly positive. Therefore, there can be no cycles in $G(\sigma_{(v)})$, that is, $\sigma_{(v)}$ is a feasible task processing order.\qed
%\begin{figure}[!t]
%\begin{minipage}{0.49\textwidth}
%\centering \includegraphics[angle=-90,trim=1cm 1cm 6cm 1cm, clip, width=0.99\textwidth]{feasibN1-1.eps}\\
%\end{minipage}
%\hfill
%\begin{minipage}{0.49\textwidth}
%\centering \includegraphics[angle=-90,trim=1cm 1cm 6cm 1cm, clip, width=0.99\textwidth]{feasibN1-2.eps}
%\end{minipage}
%\caption{\label{figFeasibility} Reversing an arc in a critical block with an alternative path}
%\end{figure}
%\end{proof}
\emph{Sketch of Proof} Suppose by contradiction that $G(\sigma_{(v)})$ has a cycle; this means that there exists an alternative path from $x$ to $y$ in $G(\sigma)$. This, together with $(x,y)$ being critical, leads to the inequality $r_x^k + p_x^k \geq r_x^k + p_x^k + p_{SJ_x}^k + p_{PJ_y}^k$ for some $k=1,2,3$, which contradicts the fact that all task durations are strictly positive. \qed

Notice that feasibility means that local search is automatically limited to the subspace of feasible task orders. It has the additional advantage of avoiding feasibility checks for the neighbours, hence increasing the efficiency of the local search procedure (reducing computational load) and avoiding the loss of feasible solutions that is usually encountered for feasibility checking procedures (cf. \cite{DellAmico1993}).

%%
%% 2) CONNECTIVITY
%%

%% LEMMA 1
%\begin{proposition}\label{propVs(s0)}
%Let $\sigma \in \Sigma$ be a feasible task processing order, $G(\sigma) =(V, A \cup E(\sigma))$ its disjunctive graph and $\sigma_0$ an optimal processing order. Let
%\begin{equation}
%V_\sigma(\sigma_0) = \{v=(x,y) \in E(\sigma) : v \text{ is critical }, (y,x) \in \overline{E(\sigma_0)} \}
%\end{equation}
%where $\overline{E(\sigma)}$ denotes the transitive closure of $E(\sigma)$; i.e., $V_\sigma(\sigma_0)$ is the set of critical arcs $(x,y)$ in $E(\sigma)$ such that there exists a path from $y$ to $x$ in $E(\sigma_0)$.
%If holds that, if $\sigma$ is not optimal, $V_\sigma(\sigma_0) \neq \emptyset$ or, equivalently, if $V_\sigma(\sigma_0) = \emptyset$ then $\sigma$ is optimal.
%\end{proposition}
%\begin{proof}
%We shall first prove that, if $\sigma$ is not optimal, there is at least a critical arc in $E(\sigma)$. Secondly, we shall prove that if $\sigma$ is not optimal, then there exists some critical arc $v=(x,y)$ such that $(y,x) \in \overline{E(\sigma_0)}$.
%
%If there are no critical arcs in $E(\sigma)$, that means that all critical arcs in $G(\sigma)$ belong to $A$. Therefore, for all $i$, all critical paths in $G^i(\sigma)$ belong to $A$. Hence, in each $G^i(\sigma)$ there exists a critical path where all arcs belong to $A$ and such path is optimal in $G^i(\sigma)$ (for every $i$, a path where all arcs belong to the same job is a lower bound of $C_{max}^i$). Therefore, $\sigma$ is optimal.
%
%Let us now asume that all critical arcs $v=(x,y) \in E(\sigma)$ verify that $(x,y) \in \overline{E(\sigma_0)}$. The set of critical arcs in $G(\sigma)$ is the union of the set of critical arcs across all parallel disjunctive graphs. Therefore, the assumption means that for all $i$ all critical arcs in $E^i(\sigma)$ belong to the transitive closure $\overline{E^i(\sigma_0)}$. Hence, a critical path $P^i$ in $G^i(\sigma)$ is also critical in $\overline{G^i(\sigma_0)} = (V, A \cup \overline{E^i(\sigma_0)})$. Let $R^i$ denote an arbitrary critical path in $\overline{G^i(\sigma_0)}$ and let $\|R^i\|$ denote its length. Since $P^i$ is critical in $\overline{G^i(\sigma_0)}$, it holds that:
%\begin{equation*}
%\forall i, C_{max}^i(\sigma) = \|P^i\| \leq \|R^i\| = C_{max}^i(\sigma_0)
%\end{equation*}
%But, since $\sigma_0$ is optimal, $E[C_{max}(\sigma_0)] \leq E[C_{max}(\sigma)]$, that is
%\begin{equation*}
%\frac {1} {4} (\|R^1\| + 2\|R^2\|+\|R^3\|) \leq \frac {1} {4} (\|P^1\| + 2\|P^2\|+\|P^3\|)
%\end{equation*}
%Both inequalities mean that $E[C_{max}(\sigma_0)] = E[C_{max}(\sigma)]$ and, hence, $\sigma$ is optimal.
%\end{proof}
%% END OF LEMMA 1

%% THEOREM: CONNECTIVITY
\begin{theorem}\label{thConnectivityH2}
$H$ verifies the connectivity property, that is, for every non-optimal task processing order $\sigma$ we may build a finite sequence of transitions of $H$ leading from $\sigma$ to a globally optimal processing order $\sigma_0$.
\end{theorem}
%% Proof
%\begin{proof}
%Let $\sigma_0$ be any optimal processing order and let the sequence $\{\lambda_k\}_{k \geq 0}$ of processing orders be given by the following recursive definition:
%\begin{gather*}
%\lambda_0 = \sigma\\
%\lambda_{k+1} \text{ is obtained from } \lambda_k \text{ by reversal of an arc } v \in V_{\lambda_k}(\sigma_0)
%\end{gather*}
%
%Notice that the reversal of an arc in $V_{\lambda_k}(\sigma_0)$ is a move from $H$ so, by Theorem~\ref{thFeasibilityH2}, $\forall k \; \lambda_k \in \Sigma$ (i.e., all task processing orders in the sequence are feasible solutions). Let us prove that the above sequence is finite. For any processing order $\sigma \in \Sigma$, we define the following sets:
%\begin{align*}
%M_{\sigma}(\sigma_0) & = \{v=(x,y) \in E(\sigma) : (y,x) \in \overline{E(\sigma_0)} \}\\
%\overline{M_{\sigma}(\sigma_0)} & = \{v=(x,y) \in \overline{E(\sigma)} : (y,x) \in \overline{E(\sigma_0)} \}
%\end{align*}
%Clearly, $V_\sigma(\sigma_0) \subset M_\sigma(\sigma_0)$ and $M_\sigma(\sigma_0) \subset \overline{M_\sigma(\sigma_0)}$. Let $\|M_\sigma(\sigma_0)\|$ and $\|\overline{M_\sigma(\sigma_0)}\|$ denote their cardinals.
%By definition of $\lambda_k$, if $\|\overline{M_{\lambda_k}(\sigma_0)}\| > 0$ then
%\begin{equation*}
%\|\overline{M_{\lambda_{k+1}}(\sigma_0)}\| = \|\overline{M_{\lambda_k}(\sigma_0)}\| -1.
%\end{equation*}
%Therefore, for $k^\star = \|\overline{M_\sigma(\sigma_0)}\|$, we have that $\|\overline{M_{\lambda_{k^\star}}(\sigma_0)}\| = 0$. Given that $V_\sigma(\sigma_0) \subseteq M_\sigma(\sigma_0) \subseteq \overline{M_\sigma(\sigma_0)}$, this implies that $V_{\lambda_{k^\star}}(\sigma_0) = \emptyset$ so, by Proposition~\ref{propVs(s0)}, $\lambda_{k^\star}$ is optimal.
%\end{proof}
%% Summary of proof
\emph{Sketch of proof.} The first step is to prove that if $V_\sigma(\sigma_0) = \{v=(x,y) \in R(\sigma) : v \text{ is critical }, (y,x) \in \overline{R(\sigma_0)} \} = \emptyset$ then $\sigma$ is optimal, where $\overline{R(\sigma_0)}$ denotes the transitive closure of $R(\sigma_0)$. To do this, first prove that if $\sigma$ is not optimal, there exists at least a critical arc in $R(\sigma)$. Then prove that if $\sigma$ is not optimal, there exists some critical arc $(x,y) \in R(\sigma)$ such that $(y,x) \in \overline{R(\sigma_0)}$, using the fact that the set of critical arcs in $G(\sigma)$ is the union of critical arcs in $G^i(\sigma)$ for all $i$. The second step is to define the sequence $\{\lambda_k\}_{k \geq 0}$, where $\lambda_0 = \sigma$, and $\lambda_{k+1}$ is obtained from $\lambda_k$ by reversal of an arc $v \in V_{\lambda_k}(\sigma_0)$. To prove that this sequence is finite, define $M_{\sigma}(\sigma_0) = \{v=(x,y) \in R(\sigma) : (y,x) \in \overline{R(\sigma_0)} \}$ and $\overline{M_{\sigma}(\sigma_0)} = \{v=(x,y) \in \overline{R(\sigma)} : (y,x) \in \overline{R(\sigma_0)} \}$, notice that if $\|\overline{M_{\lambda_k}(\sigma_0)}\| > 0$ then $\|\overline{M_{\lambda_{k+1}}(\sigma_0)}\| = \|\overline{M_{\lambda_k}(\sigma_0)}\| -1$ and use the property from step 1. \qed

As mentioned above, connectivity is an important property for any neighbourhood used in local search. It ensures the non-existence of starting points from which local search cannot reach a global optimum. It also ensures asymptotic convergence in probability to a globally optimal order. Additionally, although the neighbourhood structure is used in a heuristic procedure in this paper, the connectivity property would allow to design an exact method for fuzzy job shop, like the well-known Branch and Bound for the crisp case.

%------------------------------------------------------
% EXPERIMENTAL RESULTS
%------------------------------------------------------
%
\section{Experimental Results}\label{secExperimentalResults}
%We present results for new LS (H2) compared to old LS (H1)
% They experimentally support the theoretical results
% Son la boOOOomba :)
%% Objective
%This section presents some experimental results to illustrate the performance of both neighbourhood structures $H_1$ and $H_2$. We will use the memetic algorithm described in Section~\ref{secMA}, using both neighbourhoods in the local search.

Hybrid methods combining a genetic algorithm with local search usually improve the quality of results obtained when these methods are used independently (see for instance \cite{Ishibuchi1998}, \cite{GonzalezRodriguez2007b}, \cite{Vela2008}). The usual approach is to apply local search to every chromosome right after this chromosome has been generated. The resulting algorithm is called a \emph{memetic algorithm} (MA). To obtain experimental results for the proposed neighbourhood we shall then use a memetic algorithm proposed in \cite{GonzalezRodriguez2007b} for the fuzzy job shop, but substituting the local search procedure used therein with the local search that results from using the new neighbourhood $H$ together with simple hill-climbing.

%% 1) COMPARISON TO OTHER METHODS FROM LITERATURE: SA (FORTEMPS) AND GA (SAKAWA)
%\subsection{Comparison to Other Methods in the Literature}
The goal of the first set of experiments is to analyse the performance of the MA using the new local search compared to other methods from the literature. First, we consider a simulated annealing (SA) method from~\cite{Fortemps1997}, where the author generates fuzzy versions of well known benchmark problems~\cite{Applegate1991}: FT06 of size $6\times6$ and LA11, LA12, LA13 and LA14 of size $20 \times 5$. The only difference is that we generate 3-point fuzzy numbers or TFNs as fuzzy durations, instead of 6-point fuzzy numbers. The durations are symmetric TFNs, so the optimal solution to the crisp problem provides a lower bound for the fuzzified version (see~\cite{Fortemps1997} and \cite{GonzalezRodriguez2007b} for further detail). In total, a family of 10 instances of fuzzy job shop are generated from each original benchmark problem.

The MA has been run 30 times for each of the 50 fuzzy job shop instances with the following configuration: population size 20 and number of generations 50 (denoted 20/50 in the following). Table~\ref{tabMAvsFortemps} shows the results obtained by the MA compared to the results SA reported in~\cite{Fortemps1997}. It contains expected makespan values of the best, average and worst solution for the 10 fuzzy instances of the same crisp problem. It can be seen that the MA performs better than the SA in all cases. In fact, the best expected makespan value obtained by the MA coincides with the optimal solution for the crisp problem. Regarding CPU times, on a 3GHz Pentium IV machine, our MA takes 0.18 seconds per run for FT06 and 1.45 in average for problems LA (no CPU times are reported for the SA in the above paper).
\begin{table}[htb]
\caption{\label{tabMAvsFortemps}MA versus SA}
\centering
\begin{tabular}{c|c|ccc}
\hline\hline
Problem & Method & Best & Avg & Worst \\
\hline\hline
 FT06 & SA &  55.02 & 55.2 & 56.01 \\
  (55)& MA &  55 & 55.03 & 55.03 \\
  \hline
 LA11 & SA &  1222 & 1222 & 1222 \\
 (1222) & MA &  1222 & 1222 & 1222\\
  \hline
 LA12 & SA &  1041 & 1046.81 & 1056.35 \\
 (1039) & MA & 1039 & 1039.67 & 1041.75\\
  \hline
 LA13 & SA &  1050 & 1155.07 & 1181.76 \\
 (1150) & MA & 1150 & 1150 & 1150\\
  \hline
 LA14 & SA &  1292 & 1292 & 1292 \\
 (1292) & MA & 1292 & 1292 & 1292\\
 \hline\hline
\end{tabular}
\end{table}

Now, we compare the MA with the GA proposed in \cite{Sakawa2000}, denoted GA00, according to the implementation and parameters used in \cite{GonzalezRodriguez2008}. The problem instances are those proposed in~\cite{Sakawa2000}, and comprise three problems of size $6\times6$, denoted S6-1, S6-2 and S6-3, and three problems of size $10\times10$, denoted S10-1, S10-2 and S10-3. Table~\ref{tabMAvsGASakawa} shows the expected makespan results obtained using the same parameter setting for MA as above, 20/50, and 100/200 for GA00, both ensuring the convergence of each algorithm. It is clear that MA outperforms GA00. Moreover, the CPU time is considerably reduced, going from 33.65 seconds per run to 1.91 seconds per run in average for problems of size $10\times10$, keeping convergence curves identical, both for the average population fitness and for the fitness of the best individual.
\begin{table}[htb]
\caption{\label{tabMAvsGASakawa}MA versus GA00}
\centering
\begin{tabular}{c|c|ccc}
\hline\hline
Problem & Method & Best & Avg & Worst \\
\hline\hline
 S6-1 & GA00 & 82.25 & 83.11 & 86.25  \\
      & MA & 79.75 & 79.75 & 79.75 \\
      \hline
 S6-2 & GA00 & 75.25 & 75.69 & 78.75  \\
      & MA & 70.25 & 70.58 & 75.25 \\
       \hline
 S6-3 &  GA00 & 66.25 & 66.25 & 66.25  \\
      & MA & 66.25 & 66.25 & 66.25 \\
       \hline

 S10-1 & GA00 & 133.50 & 135.51 & 138.50 \\
       & MA & 129.00 & 129.74 & 133.00  \\
        \hline
 S10-2 & GA00 & 127.50 & 128.38 & 132.75  \\
       & MA & 123.75 & 126.67 & 127.75 \\
        \hline
 S10-3 & GA00 & 116.50 & 121.08 & 125.00 \\
       & MA & 115.00 & 115.00 & 115.00 \\
 \hline\hline
\end{tabular}
\end{table}

%% 2) COMPARISON OF LOCAL SEARCH STRATEGIES: H+HC VERSUS H'+SD (GONZALEZRODRIGUEZ2007b)
%%
%% OPTION 1: 2 LOCAL SEARCH METHODS, H+HC VERSUS H'+SD (COMBINING NEIGHBOURHOOD AND ACCEPTANCE CRITERION)
Another interesting comparison is that of the MA using the new local search with the MA from \cite{GonzalezRodriguez2007b}, denoted MA07, that uses $H'$ as neighborhood structure and steepest descent as acceptance criterion. This comparison also serves to illustrate the theoretical results showing that $H$ improves $H'$ regarding the size of sets of neighbours without worsening the expected makespan values. %
% Problem instances
The set of fuzzy problem instances are those used in \cite{GonzalezRodriguez2007b} and they are generated from well-known benchmark problems for crisp JSP: FT10 of size $10\times10$, LA24 of size $10 \times 15$, FT20 of size $20\times 5$ and ABZ7 of size $20\times 15$, following the same method as above, so the optimal solution to the crisp problem provides a lower bound for the fuzzified version. The objective is to compare both MA and MA07 on larger (and more difficult) problems. In total, a family of 10 instances of fuzzy job shop are generated from each original benchmark problem.

%% Results (1): Cmax values obtained with 2 versions of MA (with hill-climbing and H or steepest descent with H')
Table~\ref{tabMAresults} shows the average across the 10 instances of the same family of best, mean and worst solution obtained in 30 runs of each method with the parameter settings from~\cite{GonzalezRodriguez2007b}. The optimal makespan of the crisp problem is shown between brackets under the problem's name. We see that the expected makespan values are almost identical for both methods, with differences between mean relative errors w.r.t. the expected makespan lower bound about 0.1\% in average.
% TABLE with makespan results of MA
\begin{table}[ht]
\caption{\label{tabMAresults} Makespan comparison between MA and MA07.}
\begin{center}
\begin{tabular}{c|c|ccc}
\hline\hline
Problem & Method & Best & Avg & Worst \\
\hline\hline
FT10& MA07 & 934.05 & 938.29 & 952.78\\
(930) & MA & 933.85 & 939.09 & 952.98\\
\hline
FT20 & MA07 & 1165.35 & 1175.34 & 1180.90\\
(1165) & MA & 1165.35 & 1175.12 & 1181.03\\
\hline
LA24 & MA07 & 942.30 & 949.15 & 963.65\\
(935) & MA & 942.93 & 948.56 & 962.65\\
\hline
ABZ7 & MA07 & 677.90 & 686.18 & 693.38\\
(656) & MA & 677.48 & 685.77 & 693.23\\
\hline\hline
\end{tabular}
\end{center}
\end{table}

%% Results (2): Reduction of time/number of neighbours H+SH w.r.t H'+SD
%% TABLE with time and neighbours evaluated
\begin{table}[ht]
\caption{\label{tabNeighbours+Time} Total Number of Neighbours and CPU Time for MA and MA07}
\begin{center}
\begin{tabular}{c|c|cc}
\hline\hline
Problem & Method & Neigh. & CPU\\
\hline\hline
FT10 & MA07 & 2.11E+06 & 50.45\\
    & MA & 9.43E+05 & 26.94\\
\hline
FT20 & MA07 & 5.73E+06 & 125.87\\
     & MA & 2.26E+06 & 56.82\\
\hline
LA24 & MA07 & 3.53E+06 & 122.26\\
     & MA & 1.29E+06 & 52.65\\
\hline
ABZ7 & MA07 & 1.24E+07 & 852.45\\
     & MA & 3.25E+06 & 248.30\\
 \hline\hline
\end{tabular}
\end{center}
\end{table}
To further illustrate the efficiency gain of the new local search, combining neighbourhood $H$ with simple hill climbing, Table~\ref{tabNeighbours+Time} reports the total number of evaluated neighbours and the time required by each method, MA and MA07. The CPU times, on a 3GHz Pentium IV machine, refer to the seconds taken per run of the MA in average. There is a clear advantage in using the new local search. It is remarkable that the new neighbourhood $H$ combined with the new selection criterion obtains identical makespan values to MA07, whilst reducing the number of evaluated neighbours more than 63\% in average. This agrees with the theoretical properties regarding $H$ and $H'$. The reduction of evaluated neighbours is reflected in a reduction of CPU times, despite it being more expensive to compute the neighbourhood $H$ than $H'$ because the search for critical arcs is performed through tree solution graphs instead of one. The average time reduction of MA w.r.t. MA07 is as much as 57\%, with the remarkable case of the largest problem ABZ7, with a time reduction superior to 70\%. Notice that the reduction in time is increased with the size of the problem, something natural since evaluating new neighbours has a higher computational cost in larger problems.

We have also conducted experiments (not reported here) to compare the gain using $H$ instead $H'$ independently from the acceptance criterion. The obtained makespan values equal those in Table~\ref{tabMAresults}, with an average reduction in the number of neighbours near 20\%, which again increases with the problem size. These results are the same regardless of the acceptance criterion used, steepest descent and simple hill-climbing.

%% 3) SYNERGY
%\subsection{Synergy}
The objective of the final set of experiments is to illustrate the synergy gained when combining the GA and the local search procedure (LS hereafter), by evaluating LS and GA separately, as well as both of them combined in the MA. This experiment also serves to test that there is no premature convergence to local optima. For results to be comparable, it is necessary to test all methods under equivalent conditions regarding computational resources.
\begin{table}[htb]
\caption{\label{tabGA-AvsGA-LSA} Synergy: Comparison between LS, GA and MA.}
\centering
\begin{tabular}{{c|c|ccc}}
\hline\hline
Problem & Method & Best & Avg & Worst \\
\hline\hline
     & LS  & 965.50 & 989.13 & 1005.92 \\
FT10 & GA  & 937.40 & 956.19 & 981.45 \\
(930) & MA & 933.85 & 939.09 & 952.98 \\
\hline
     & LS  & 985.70 & 1004.70 & 1017.20 \\
LA24 & GA  & 960.80 & 981.53 & 1005.50 \\
(935) & MA & 942.93 & 948.56 & 962.65 \\
\hline
      & LS   & 1266.83 & 1298.85 & 1322.78\\
FT20  & GA   & 1178.55 & 1188.36 & 1208.80\\
(1165) & MA & 1165.35 & 1175.12 & 1181.03 \\
\hline
    & LS   & 724.00 & 734.59 & 742.13 \\
ABZ7 & GA   & 690.75 & 706.31 & 720.90 \\
(655) & MA & 677.48 & 685.77 & 693.23 \\
\hline\hline
\end{tabular}
\end{table}
To do so, assuming a uniform computational load in all generations, for each problem we consider the ratio between CPU times of both GA and MA when they were executed with the same configuration (100/200) and then readjust the number of generations for the GA to obtain similar CPU times. To study LS independently, the LS algorithm is run 30 times for each instance starting from a set of 20000 random solutions, so as to evaluate the same number of individuals that the GA and obtain similar running times (about 27 seconds for FT10, 57 for FT20, 53 for LA24 and 248 for ABZ7). Table~\ref{tabGA-AvsGA-LSA} shows the expected makespan results for each family of problems. It is worth noting that for MA, the CPU time dedicated to local search is 91-96\% of the total, which can be explained by the fact that LS is applied to every chromosome in the population. We may conclude that, with the same computational time, the MA benefits from the synergy between the GA and LS and always yields better solutions than the GA or LS alone.

%------------------------------------------------------
% CONCLUSIONS
%------------------------------------------------------
%
\section{Conclusions}\label{secConclusions}
We have considered a job shop problem with uncertain processing times modelled as triangular fuzzy numbers, where the objective is to minimise the expected makespan value. We have proposed a new definition of criticality for this problem and a consequent new neighbourhood structure for local search. We have shown that the new structure improves previous proposals from the literature, reducing the set of neighbours by discarding non-improving ones. Additionally, we have shown that the new neighbourhood structure has two highly desirable properties: feasibility and connectivity, which further improve the efficiency of the resulting local search, as well as ensuring asymptotic convergence in probability to a globally optimal solution. The experimental results support the theoretical results and illustrate the potential of the new local search when used in combination with a genetic algorithm, comparing favourably with previous approaches from the literature.

In the future, makespan estimation procedures will be developed to reduce the computational cost of local search. These estimates may also be used to design a branch and bound type algorithm, taking advantage of the connectivity property. Additionally, metaheuristics other than the memetic algorithm will be proposed using the new neighbourhood, for instance, a taboo search procedure. Finally, the parallel solution graph model and consequent definitions of criticality and neighbourhood, proposed for durations represented as TFNs, may be extended to piece-wise linear fuzzy intervals, using $\alpha$-cuts.

%------------------------------------------------------
% ACKNOWLEDGEMENTS
%------------------------------------------------------

\section*{ Acknowledgments}
All authors are supported by MEC-FEDER Grant TIN2007-67466-C02-01.

%------------------------------------------------------
% REFERENCES
%------------------------------------------------------

%\bibliographystyle{aaai}
%\bibliography{../../../a2zScheduling}

\begin{thebibliography}{}

\bibitem[\protect\citeauthoryear{Applegate \& Cook}{1991}]{Applegate1991}
Applegate, D., and Cook, W.
\newblock 1991.
\newblock A computational study of the job-shop scheduling problem.
\newblock {\em ORSA Journal of Computing} 3:149--156.

\bibitem[\protect\citeauthoryear{Bierwirth}{1995}]{Bierwirth1995}
Bierwirth, C.
\newblock 1995.
\newblock A generalized permutation approach to jobshop scheduling with genetic
  algorithms.
\newblock {\em {OR} Spectrum} 17:87--92.

\bibitem[\protect\citeauthoryear{Brucker \& Knust}{2006}]{Brucker2006}
Brucker, P., and Knust, S.
\newblock 2006.
\newblock {\em Complex Scheduling}.
\newblock Springer.

\bibitem[\protect\citeauthoryear{Celano, Costa, \& Fichera}{2003}]{Celano2003}
Celano, G.; Costa, A.; and Fichera, S.
\newblock 2003.
\newblock An evolutionary algorithm for pure fuzzy flowshop scheduling
  problems.
\newblock {\em International Journal of Uncertainty, Fuzziness and
  Knowledge-Based Systems} 11:655--669.

\bibitem[\protect\citeauthoryear{Dell'~Amico \& Trubian}{1993}]{DellAmico1993}
Dell'~Amico, M., and Trubian, M.
\newblock 1993.
\newblock Applying tabu search to the job-shop scheduling problem.
\newblock {\em Annals of Operational Research} 41:231--252.

\bibitem[\protect\citeauthoryear{Dubois \& Prade}{1986}]{Dubois1988}
Dubois, D., and Prade, H.
\newblock 1986.
\newblock {\em Possibility Theory: An Approach to Computerized Processing of
  Uncertainty}.
\newblock New York (USA): Plenum Press.

\bibitem[\protect\citeauthoryear{Dubois, Fargier, \&
  Fortemps}{2003}]{Dubois2003}
Dubois, D.; Fargier, H.; and Fortemps, P.
\newblock 2003.
\newblock Fuzzy scheduling: Modelling flexible constraints vs. coping with
  incomplete knowledge.
\newblock {\em European Journal of Operational Research} 147:231--252.

\bibitem[\protect\citeauthoryear{Fortemps}{1997}]{Fortemps1997}
Fortemps, P.
\newblock 1997.
\newblock Jobshop scheduling with imprecise durations: a fuzzy approach.
\newblock {\em IEEE Transactions of Fuzzy Systems} 7:557--569.

\bibitem[\protect\citeauthoryear{Gonz\'alez~Rodr\'{\i}guez \bgroup \em et
  al.\egroup }{2008}]{GonzalezRodriguez2008}
Gonz\'alez~Rodr\'{\i}guez, I.; Puente, J.; Vela, C.~R.; and Varela, R.
\newblock 2008.
\newblock Semantics of schedules for the fuzzy job shop problem.
\newblock {\em IEEE Transactions on Systems, Man and Cybernetics, Part A}
  38(3):655--666.

\bibitem[\protect\citeauthoryear{Gonz\'alez~Rodr\'{\i}guez, Vela, \&
  Puente}{2007}]{GonzalezRodriguez2007b}
Gonz\'alez~Rodr\'{\i}guez, I.; Vela, C.~R.; and Puente, J.
\newblock 2007.
\newblock A memetic approach to fuzzy job shop based on expectation model.
\newblock In {\em Proceedings of IEEE International Conference on Fuzzy
  Systems},  692--697.
\newblock London: IEEE.

\bibitem[\protect\citeauthoryear{Ishibuchi \& Murata}{1998}]{Ishibuchi1998}
Ishibuchi, H., and Murata, T.
\newblock 1998.
\newblock A multi-objective genetic local search algorithm and its application
  to flowshop scheduling.
\newblock {\em IEEE Transactions on Systems, Man, and Cybernetics--Part C} 67(3):392--403.

\bibitem[\protect\citeauthoryear{Liu \& Liu}{2002}]{Liu2002}
Liu, B., and Liu, Y.~K.
\newblock 2002.
\newblock Expected value of fuzzy variable and fuzzy expected value models.
\newblock {\em IEEE Transactions on Fuzzy Systems} 10:445--450.

\bibitem[\protect\citeauthoryear{Sakawa \& Kubota}{2000}]{Sakawa2000}
Sakawa, M., and Kubota, R.
\newblock 2000.
\newblock Fuzzy programming for multiobjective job shop scheduling with fuzzy
  processing time and fuzzy duedate through genetic algorithms.
\newblock {\em European Journal of Operational Research} 120:393--407.

\bibitem[\protect\citeauthoryear{S{\l}owi\'nski \& Hapke}{2000}]{Slowinski2000}
S{\l}owi\'nski, R., and Hapke, M., eds.
\newblock 2000.
\newblock {\em Scheduling Under Fuzziness}, volume~37 of {\em Studies in
  Fuzziness and Soft Computing}.
\newblock Physica-Verlag.

\bibitem[\protect\citeauthoryear{Tavakkoli-Moghaddam, Safei, \&
  Kah}{2008}]{Tavakkoli-Moghaddam2008}
Tavakkoli-Moghaddam, R.; Safei, N.; and Kah, M.
\newblock 2008.
\newblock Accessing feasible space in a generalized job shop scheduling problem
  with the fuzzy processing times: a fuzzy-neural approach.
\newblock {\em Journal of the Operational Research Society} 59:431--442.

\bibitem[\protect\citeauthoryear{{Van Laarhoven}, Aarts, \&
  Lenstra}{1992}]{VanLaarhoven1992}
{Van Laarhoven}, P.; Aarts, E.; and Lenstra, K.
\newblock 1992.
\newblock Job shop scheduling by simulated annealing.
\newblock {\em Operations Research} 40:113--125.

\bibitem[\protect\citeauthoryear{Vela, Varela, \& Gonz\'alez}{2008}]{Vela2008}
Vela, C.~R.; Varela, R.; and Gonz\'alez, M.~A.
\newblock 2008.
\newblock Local search and genetic algorithm for the job shop scheduling
  problem with sequence dependent setup times.
\newblock {\em Journal of Heuristics} Forthcoming.

\end{thebibliography}


\end{document}
